\section{Introduction}\label{sec:intro}

Whether, and when, hostility expressed toward Israel reflects hostility
toward Jews is among the most consequential and contested questions in
public life today. The answer shapes campus policy, law and community
safety. In 2017 the European Parliament called on the Member States and the
institutions of the European Union to adopt and apply the working definition
of antisemitism of the International Holocaust Remembrance Alliance (IHRA)
\citep{ep2017antisemitism}; in October 2024 the European Commission reported
that 25 Member States had adopted or endorsed it \citep{ec2024antisemitism}.
In the United States, a 2019 executive order directs the agencies that
enforce Title~VI of the Civil Rights Act, which prohibits discrimination on
the basis of race, color or national origin in federally funded programmes,
to consider the IHRA definition and, where they are useful as evidence of
discriminatory intent, its contemporary examples. The order also bars the
agencies from diminishing or infringing any right protected by federal law
or the First Amendment, and its statement of policy refers to harassment in
schools and on university and college campuses \citep{eo13899}; a 2025 order
reaffirmed it \citep{eo14188}. Several of the IHRA examples concern Israel
\citep{ihra2016}. Scholars who regard the IHRA text as unclear in key
respects have proposed the Jerusalem Declaration on Antisemitism as an
alternative or corrective \citep{jda2021}, and the Nexus Document offers a
third formulation \citep{nexus2021}. The three texts draw the line around
Israel-related speech in different places (\cref{sec:bg-definitions}). The
same line runs through incident statistics. The incident lists studied in
the first paper of this series differ most in how they treat Israel- or
Zionism-related expression \citep{kornsteinlab2026darkfigure}, so counts of
incidents, and the security planning that rests on them, inherit a choice of
definition.

\paragraph{Normative and empirical claims}
A definition is a normative instrument. It says which speech and conduct
ought to be called antisemitic, by whom and with what consequences, and no
survey can settle that. But the competing definitions also embed empirical
claims alongside the normative ones. Take a particular statement about
Israel. One claim is that endorsing it typically expresses the same hostility
toward Jews that drives endorsement of classic antisemitic stereotypes. The
competing claim is that, in the population at large, endorsing it mainly
expresses a distinct political attitude toward a state and its policies.
Each is a claim about how attitudes are structured in a population,
statement by statement. Either can be right for some statements and wrong
for others. And each implies restrictions on a measurement model, so
measurement can test it. This paper builds the tools for that test. It does
not choose between the definitions, and it is designed so that its results
would be equally usable, and equally open to checking, by people on every
side of the debate.

\IfFileExists{figures/fig-bifactor.tex}{% =====================================================================
%  Figure: the bifactor measurement model and its comparison model (LEAD).
%  Path diagram. Loadings are SCHEMATIC (illustrative values, not data,
%  not the simulation design): arrow widths grow with the slope, and the bar
%  in each item box is ECV_j = a_jG^2 / (a_jG^2 + a_js^2), computed in TikZ
%  from the same values.  Items 1-5 are stereotype items (set S), items 6-11
%  Israel-related items (set I); items 1, 2, 6, 7 are anchors (set A).
%  Requires the house colours defined in preamble.tex.
% =====================================================================
\begin{figure}[t]
\centering
\begin{tikzpicture}[x=1cm,y=1cm,
    item/.style={rectangle,rounded corners=1.6pt,draw=LabNavyLight,line width=0.6pt,
                 fill=white,minimum width=0.74cm,minimum height=0.62cm,inner sep=0pt},
    anchoritem/.style={item,draw=LabAccent,line width=1.0pt,fill=LabSand},
    gfac/.style={circle,fill=LabNavy,draw=LabNavy,text=white,minimum size=0.98cm,
                 inner sep=0pt,font=\normalsize},
    sfac/.style={circle,fill=white,draw=LabNavyMid,line width=0.9pt,text=LabNavyMid,
                 minimum size=0.86cm,inner sep=0pt,font=\normalsize},
    gl/.style={LabNavy,-{Stealth[length=3.4pt,width=2.8pt]},shorten >=0.5pt},
    sl/.style={LabNavyMid,-{Stealth[length=3.4pt,width=2.8pt]},shorten >=0.5pt},
    panel/.style={font=\footnotesize\bfseries,text=LabNavy,anchor=south west,inner sep=0pt},
    note/.style={font=\scriptsize,text=LabGrey,inner sep=1pt,align=center},
    leg/.style={font=\scriptsize,text=LabInk,anchor=west,inner sep=1pt},
    pl/.style={font=\scriptsize,text=#1,fill=white,inner sep=0.8pt,fill opacity=0.92,text opacity=1},
  ]
  % ---------------- schematic item parameters ----------------
  % index / x position / a_G / a_s / specific factor (S or I) / anchor (1/0)
  \def\items{1/0.00/1.30/0.60/S/1, 2/0.92/1.20/0.75/S/1, 3/1.84/1.10/0.50/S/0,
             4/2.76/1.25/0.80/S/0, 5/3.68/1.00/0.65/S/0,
             6/4.98/1.00/0.70/I/1, 7/5.90/0.80/0.95/I/1, 8/6.82/0.50/1.20/I/0,
             9/7.74/0.90/0.85/I/0, 10/8.66/0.35/1.30/I/0, 11/9.58/1.10/0.60/I/0}
  \def\ylev{2.40}\def\yspec{-2.10}
  \def\barw{0.58}

  % =========================== panel (a) ===========================
  \node[panel] at (-0.42,3.22) {(a)\enspace Bifactor model
    \textnormal{\textcolor{LabGrey}{(primary)}}};

  % factors
  \node[gfac] (G) at (4.79,\ylev) {$G$};
  \node[sfac] (S) at (1.84,\yspec) {$S$};
  \node[sfac] (I) at (7.28,\yspec) {$I$};

  % items: box, label, ECV bar
  \foreach \j/\x/\ag/\as/\f/\an in \items {
    \ifnum\an=1 \node[anchoritem] (y\j) at (\x,0) {};
    \else \node[item] (y\j) at (\x,0) {};\fi
    \node[font=\footnotesize,text=LabInk,inner sep=0pt] at (\x,0.085) {$Y_{\j}$};
    \pgfmathsetmacro{\ecv}{\ag*\ag/(\ag*\ag+\as*\as)}
    \fill[LabNavyLight!55] (\x-\barw/2,-0.215) rectangle (\x+\barw/2,-0.15);
    \fill[LabNavy] (\x-\barw/2,-0.215) rectangle (\x-\barw/2+\ecv*\barw,-0.15);
  }
  % loadings (drawn on the background so that boxes sit on top)
  \begin{scope}[on background layer]
    \foreach \j/\x/\ag/\as/\f/\an in \items {
      \pgfmathsetmacro{\wg}{0.22+0.42*\ag}
      \pgfmathsetmacro{\ws}{0.22+0.42*\as}
      \draw[gl,line width=\wg pt] (G) -- (y\j.north);
      \draw[sl,line width=\ws pt] (\f) -- (y\j.south);
    }
  \end{scope}

  % loading labels on the highlighted item j = 9
  \node[pl=LabNavy] at ($(G)!0.8!(y9.north)$) {$a_{jG}$};
  \node[pl=LabNavyMid] at ($(I)!0.55!(y9.south)$) {$a_{jI}$};

  % the item-level estimand, for the highlighted item
  \node[font=\footnotesize,text=LabNavy,anchor=west,inner sep=1pt] (ecv) at (6.3,3.0)
    {$\mathrm{ECV}_j=\dfrac{a_{jG}^2}{a_{jG}^2+a_{jI}^2}$};
  \node[note,anchor=north west,align=left] (ecvn) at (6.37,2.5)
    {share of item $j$'s common variance\\[-1pt]that is due to $G$ (here $j=9$)};

  % factor labels
  \node[note,anchor=east,align=right] at ([xshift=-4pt]G.west)
    {\textcolor{LabNavy}{general factor}\\[-1pt]loads on every item};
  \node[note,anchor=north] at ([yshift=-2pt]S.south)
    {\textcolor{LabNavyMid}{specific factor}: stereotype items $\mathcal{S}$};
  \node[note,anchor=north] at ([yshift=-2pt]I.south)
    {\textcolor{LabNavyMid}{specific factor}: Israel-related items $\mathcal{I}$};

  % legend
  \begin{scope}[shift={(-0.37,-3.32)}]
    \node[anchoritem,minimum width=0.42cm,minimum height=0.3cm] (la) at (0,0) {};
    \node[leg] (lat) at (0.28,0) {anchor item ($j\in\mathcal{A}$): assumed invariant across groups};
    \draw[LabNavyLight!55,line width=2.3pt] (6.95,0) -- (7.42,0);
    \draw[LabNavy,line width=2.3pt] (6.95,0) -- (7.22,0);
    \node[leg] at (7.52,0) {bar: $\mathrm{ECV}_j$ (schematic)};
  \end{scope}

  % ---------------- separator ----------------
  \draw[LabRule,line width=0.4pt] (10.42,-3.45) -- (10.42,3.35);

  % =========================== panel (b) ===========================
  \begin{scope}[shift={(10.9,0)}]
    \node[panel] at (-0.15,3.22) {(b)\enspace Correlated factors
      \textnormal{\textcolor{LabGrey}{(comparison)}}};
    \node[sfac] (S2) at (0.76,\yspec) {$\zeta_S$};
    \node[sfac] (I2) at (3.00,\yspec) {$\zeta_I$};
    \foreach \j/\x/\f/\an in {1/0.00/S/1,2/0.38/S/1,3/0.76/S/0,4/1.14/S/0,5/1.52/S/0,
                              6/2.05/I/1,7/2.43/I/1,8/2.81/I/0,9/3.19/I/0,10/3.57/I/0,11/3.95/I/0}{
      \ifnum\an=1 \node[anchoritem,minimum width=0.27cm] (z\j) at (\x,0) {};
      \else \node[item,minimum width=0.27cm] (z\j) at (\x,0) {};\fi
      \begin{scope}[on background layer]
        \draw[sl,line width=0.5pt] (\f2) -- (z\j.south);
      \end{scope}
    }
    \draw[LabNavyMid,line width=0.8pt,{Stealth[length=3.4pt,width=2.8pt]}-{Stealth[length=3.4pt,width=2.8pt]},
          shorten <=1pt,shorten >=1pt]
      (S2.south) to[out=-55,in=-125] node[pos=0.5,below=0.5pt,font=\footnotesize,text=LabNavyMid]
      {$\rho_{SI}$} (I2.south);
    \node[note,anchor=north,text width=4.2cm] at (1.97,1.95)
      {no general factor: the same items\\[-1pt]load on two correlated factors};
    \node[note,anchor=north,text width=4.3cm] at (1.97,-3.55)
      {latent correlation between the\\[-1pt]stereotype and Israel-attitude factors};
  \end{scope}
\end{tikzpicture}
\caption{\textbf{The measurement model.} \emph{(a)} The bifactor graded
response model of \cref{sec:model}. Every item loads on the general factor
$G$, which carries the variance that all items share. Each item also loads on
exactly one specific factor: $S$ for the classic-stereotype items
$\mathcal{S}$ and $I$ for the Israel-related items $\mathcal{I}$. Responses are
ordinal, with $\Pr(Y_{ij}\ge c\mid\bm\theta_i)
=F(a_{jG}\theta_{iG}+a_{js}\theta_{is}-b_{jc})$ for item $j$'s specific factor
$s$, and the three factors are uncorrelated (\cref{eq:model-latent}). The
boundary estimand of \cref{def:boundary} is the share of an item's common
variance that is due to $G$, $\mathrm{ECV}_j$ (shown for item~9); the bar in
each box shows it, and arrow widths grow with the slopes. Anchor items
(orange outline, $\mathcal{A}$) are assumed to function identically in every
group and carry group comparisons when other items show differential item
functioning (\cref{prop:anchor}). \emph{(b)} The pre-registered comparison
model: the same items, no general factor, and two correlated domain factors
$\zeta_S$ and $\zeta_I$, whose correlation $\rho_{SI}$ is the second boundary
estimand. Every such model is also a bifactor model with loadings
proportional within each domain, so it identifies $\rho_{SI}$ but not the
item-level shares, which can lie anywhere in $[\rho_{SI}^2,1]$
(\cref{prop:model-sl}). Schematic: the loadings are illustrative values, not
estimates or simulation inputs.}
\label{fig:bifactor}
\end{figure}
}{}%
\paragraph{A latent-variable map}
We state the empirical question in a bifactor graded response model
(\cref{fig:bifactor}). Respondents rate two kinds of statements: classic
stereotypes about Jews, and statements about Israel, its policies and its
supporters. A general factor $G$ loads on every item. Two specific factors
absorb what is left: $S$, which the stereotype items share only with each
other, and $I$, which the Israel-related items share only with each other.
For each Israel-related item, the share of its common variance that is due to
$G$, an item-level explained common variance $\mathrm{ECV}_j$, places the
item on a map. Near $0$, the item's common variance is what it shares with
the other Israel-related items alone; near $1$, it is what it shares with the
classic stereotypes. The \emph{boundary set} $\mathcal{B}_\tau$ collects the items
whose share reaches a threshold $\tau$ fixed before any data are seen, and
sampling uncertainty sorts each item into \emph{in}, \emph{out} or
\emph{undetermined} (\cref{fig:boundary}). The map is drawn item by item,
and the same rule applies to every item. Earlier studies suggest that the
relation between the two domains is not uniform. \Citet{beattie2017}
reports that some criticisms of the Israeli government are more reliable
indicators of antisemitic sentiment than others, and \citet{staetsky2017}
that more intense anti-Israel attitudes are more likely to be accompanied by
antisemitic ones. A correlated-factors
model, with no general factor, supplies a second, coarser estimand: the
latent correlation $\rho_{SI}$ between the stereotype and the
Israel-attitude factors. On its own, that correlation cannot tell a common
latent cause from distinct but correlated attitudes: every correlated-factors
model is also a bifactor model in which the Israel-related items' share of
general-factor variance can be anywhere between $\rho_{SI}^2$ and one
(\cref{prop:model-sl}).
The general factor has no meaning of its own. It is defined by the items
that load on it, which is why the item bank and its balance across the
disputed categories are fixed in advance (\cref{sec:design-items}).

\IfFileExists{figures/fig-boundary.tex}{% =====================================================================
%  Figure: a schematic boundary map (LEAD).
%  SCHEMATIC, NOT DATA. Ten generic Israel-related items with MADE-UP point
%  estimates and 95% intervals for ECV_j (no item content, no estimates, no
%  simulation output). The classification is computed in TikZ from the
%  made-up intervals with the rule of the pre-registered plan
%  (sec:design-estimands): in B_tau if lo >= tau; out if hi < tau;
%  undetermined otherwise. Intervals are read as simultaneous (Bonferroni) 95% intervals. Primary tau = 0.5; sensitivity tau = 0.3, 0.7.
%  Requires the house colours defined in preamble.tex.
% =====================================================================
\begin{figure}[t]
\centering
\begin{tikzpicture}[x=0.86cm,y=0.5cm,
    lab/.style={font=\scriptsize,inner sep=1pt},
    tick/.style={font=\scriptsize,text=LabGrey,inner sep=1pt},
    hdr/.style={font=\scriptsize\bfseries,text=LabNavy,inner sep=1pt},
  ]
  \def\sx{10}   % x units per unit of ECV
  \def\tauz{0.5}
  % row / point estimate / lower / upper  (made up; ordered by estimate)
  \def\rows{1/0.86/0.78/0.92, 2/0.78/0.66/0.88, 3/0.71/0.55/0.84,
            4/0.63/0.44/0.79, 5/0.55/0.41/0.69, 6/0.47/0.30/0.64,
            7/0.38/0.26/0.52, 8/0.31/0.19/0.44, 9/0.22/0.12/0.35,
            10/0.12/0.05/0.22}
  \def\ytop{10.0}\def\ybot{0.0}

  % ---------------- regions ----------------
  \fill[LabTint] (\tauz*\sx,\ybot) rectangle (\sx,\ytop);
  \draw[LabRule,line width=0.4pt] (0,\ybot) rectangle (\sx,\ytop);
  \node[lab,text=LabNavyMid,anchor=south,align=center] at (0.25*\sx,\ytop+1.15)
    {mainly the Israel-specific factor $I$\\[-1pt]$\mathrm{ECV}_j<\tau$};
  \node[lab,text=LabNavy,anchor=south,align=center] at (0.75*\sx,\ytop+1.15)
    {mainly the general factor $G$\\[-1pt]$\mathrm{ECV}_j\ge\tau$};

  % ---------------- thresholds ----------------
  \foreach \t in {0.3,0.7}{
    \draw[LabAccent!55,line width=0.5pt,dash pattern=on 2pt off 1.5pt] (\t*\sx,\ybot) -- (\t*\sx,\ytop);}
  \draw[LabAccent,line width=0.9pt] (\tauz*\sx,\ybot) -- (\tauz*\sx,\ytop+0.35);
  \node[lab,text=LabAccent,anchor=south,font=\scriptsize\bfseries] at (\tauz*\sx,\ytop+0.3)
    {pre-registered threshold $\bm{\tau=0.5}$};

  % ---------------- decision rule (empty top-left corner) ----------------
  \node[lab,text=LabGrey,anchor=north west,align=left,fill=white,inner sep=2.5pt]
    at (0.1,\ytop-0.12)
    {\textcolor{LabNavy}{in $\mathcal{B}_\tau$}: lower 95\% limit $\ge\tau$\\
     \textcolor{LabNavy}{out}: upper 95\% limit $<\tau$\\
     undetermined: interval contains $\tau$\\
     \textcolor{LabAccent!75}{dashed}: sensitivity values of $\tau$};
  % ---------------- schematic stamp (empty bottom-right corner) ----------------
  \node[font=\small\itshape,text=LabGrey,anchor=south east,inner sep=2pt]
    at (\sx-0.1,\ybot+0.25) {schematic, not data};

  % ---------------- column headers ----------------
  \node[hdr,anchor=south east,align=right] at (-0.25,\ytop+0.3) {Israel-related\\[-1pt]items};
  \node[hdr,anchor=south west,align=left] at (\sx+0.3,\ytop+0.3) {classification\\[-1pt]at $\tau=0.5$};
  \node[hdr,anchor=south,align=center,text=LabGrey] at (\sx+3.55,\ytop+0.3)
    {sensitivity\\[-1pt]$\tau=0.3$\enspace$0.7$};

  % ---------------- items ----------------
  \foreach \r/\e/\lo/\hi in \rows {
    \pgfmathsetmacro{\yy}{\ytop-\r+0.5}
    \pgfmathtruncatemacro{\cls}{\lo>=\tauz ? 1 : (\hi<\tauz ? -1 : 0)}
    \ifnum\cls=0
      \def\ccol{LabGrey}\def\cfill{white}\def\ctext{undetermined}
    \else
      \def\ccol{LabNavy}\def\cfill{LabNavy}%
      \ifnum\cls=1 \def\ctext{in $\mathcal{B}_\tau$}\else\def\ctext{out}\fi
    \fi
    % interval, end caps, point
    \draw[\ccol,line width=1.1pt] (\lo*\sx,\yy) -- (\hi*\sx,\yy);
    \draw[\ccol,line width=0.6pt] (\lo*\sx,\yy-0.17) -- (\lo*\sx,\yy+0.17)
                                  (\hi*\sx,\yy-0.17) -- (\hi*\sx,\yy+0.17);
    \fill[\cfill,draw=\ccol,line width=0.8pt] (\e*\sx,\yy) circle (2.3pt);
    % row label and primary classification
    \node[tick,anchor=east] at (-0.25,\yy) {item \r};
    \node[lab,text=\ccol,anchor=west] at (\sx+0.3,\yy) {\ctext};
    % sensitivity classifications
    \foreach \t/\cx in {0.3/2.95,0.7/4.15}{
      \pgfmathtruncatemacro{\cs}{\lo>=\t ? 1 : (\hi<\t ? -1 : 0)}
      \ifnum\cs=1 \node[lab,text=LabNavy] at (\sx+\cx,\yy) {in};\fi
      \ifnum\cs=-1 \node[lab,text=LabNavy] at (\sx+\cx,\yy) {out};\fi
      \ifnum\cs=0 \node[lab,text=LabGrey] at (\sx+\cx,\yy) {und.};\fi
    }
  }
  \draw[LabRule,line width=0.4pt] (\sx+2.45,\ybot) -- (\sx+2.45,\ytop);

  % ---------------- axis ----------------
  \draw[LabNavy,line width=0.5pt] (0,\ybot) -- (\sx,\ybot);
  \foreach \v/\c in {0/LabGrey,0.3/LabAccent!80,0.5/LabAccent,0.7/LabAccent!80,1/LabGrey}{
    \draw[LabNavy,line width=0.5pt] (\v*\sx,\ybot) -- (\v*\sx,\ybot-0.25);
    \node[tick,text=\c,anchor=north] at (\v*\sx,\ybot-0.3) {$\v$};}
  \node[lab,text=LabNavy,anchor=north] at (0.5*\sx,\ybot-1.15)
    {$\mathrm{ECV}_j$: share of item $j$'s common variance due to $G$, with simultaneous 95\% interval};
  \node[lab,text=LabGrey,anchor=north west,font=\scriptsize\itshape] at (0,\ybot-2.0)
    {$\leftarrow$ mainly shared with other Israel-related items};
  \node[lab,text=LabGrey,anchor=north east,font=\scriptsize\itshape] at (\sx,\ybot-2.75)
    {mainly shared with the stereotype items $\rightarrow$};
\end{tikzpicture}
\caption{\textbf{A boundary map (schematic, not data).} What the primary
pre-registered estimands E1 and E2 (\cref{sec:design-estimands})
would look like when reported. Each row is one Israel-related item, placed by
its item-level explained common variance $\mathrm{ECV}_j$
(\cref{def:boundary}): the share of the item's common variance that it shares,
through the general factor $G$, with the classic-stereotype items, rather than
with the other Israel-related items alone, through the specific factor $I$
(\cref{fig:bifactor}). Bars are simultaneous 95\% intervals (\cref{prop:id-boundary-test}(b)). An item is \emph{in} the
boundary set $\mathcal{B}_\tau$ when its whole interval lies at or above the
pre-registered threshold $\tau=0.5$, \emph{out} when it lies wholly below,
and \emph{undetermined} when it contains $\tau$; the same rule is applied to
every item, whatever its content. The right-hand columns repeat the
classification at the sensitivity values $\tau=0.3$ and $0.7$. Positions and
intervals are made up for illustration; the items are generic and carry no
content, and nothing here is an estimate or a simulation result. The map
describes response patterns in a surveyed population. It does not classify
any utterance, speaker or context (\cref{sec:ethics-individual}).}
\label{fig:boundary}
\end{figure}
}{}%
\paragraph{Four obstacles}
Fitting such a model is easy; reading it correctly is not. We formalise four
obstacles.
\begin{enumerate}[(i),leftmargin=2.2em]
  \item \emph{Identification.} A bifactor model can be fitted to data that do
  not identify it. The decomposition is identified by heterogeneity, not by
  correlation: if, within each domain, all items but at most one have the
  same general-factor share, a whole curve of parameter values fits the
  data equally well and
  $\mathrm{ECV}_j$ is not determined, and with only two Israel-related items
  it is only bounded (\cref{thm:ident}).
  \item \emph{Invariance.} Respondents who differ in ideology, age or country
  may read an Israel-related statement differently, so that the item
  functions differently across groups. Ignoring such differential item
  functioning (DIF) biases group comparisons of latent means and of the
  boundary estimands. To first order the bias is an information-weighted
  projection of the DIF onto the loadings, and uniform DIF proportional to
  the general loadings is indistinguishable from a true difference in the
  general factor (\cref{thm:dif-bias}). Invariant anchor items restore the
  comparison, but DIF in the anchors themselves shifts every comparison
  they carry, undetectably when there are only three of them
  (\cref{prop:anchor}).
  \item \emph{Social desirability.} Direct questions on these topics may be
  answered strategically, and not necessarily in the same direction in every
  group. List experiments protect individual answers, at a variance cost
  that we quantify (\cref{prop:list-var}). When strategic answering acts as a
  uniform shift of an item's thresholds, it biases prevalences but not the
  boundary estimands (\cref{prop:list-irt}); latent-mean comparisons are
  affected only if the shift differs between groups, in which case it is a
  form of DIF.
  \item \emph{Aggregation.} Policy questions are often local, while samples
  are national. Small-area estimates must be built from a model, and scoring
  respondents first and averaging the scores afterwards shrinks every
  difference between areas (\cref{prop:eap}). Indices that report the share
  of respondents agreeing with at least $h$ of $M$ stereotype statements
  rank groups correctly under invariance and stochastic dominance, but not
  in general (\cref{prop:index}).
\end{enumerate}

\paragraph{What this paper does, and does not, do}
This is a methods paper. We set out the measurement model and its boundary
estimands, prove when they are identified and what goes wrong when the
model's assumptions fail, develop the corrections for social desirability and
the small-area estimators, evaluate all of them in a simulation study on
synthetic survey data with known ground truth, and pre-register a survey
design and analysis plan. We report \emph{no} empirical estimates. We have
fielded no survey and fitted no model to real data, and nothing here should
be read as a finding that any attitude toward Israel is, or is not,
antisemitic, or about the attitudes of any group. Survey results published by
others are quoted only as their authors report them. Notation is collected in
\cref{tab:notation}.

\IfFileExists{tables/tab-notation.tex}{% =====================================================================
%  Table: notation (LEAD).  Symbols follow SPEC.md and Sections 3--6.
% =====================================================================
\begin{table}[t]
\caption{\textbf{Notation.} Symbols used throughout. Item parameters carry
a superscript $(g)$ when they are allowed to differ across groups.}
\label{tab:notation}
\small
\renewcommand{\arraystretch}{1.08}
\newcommand{\nhead}[1]{\multicolumn{2}{@{}l}{\footnotesize\scshape\bfseries\color{LabNavy}#1}\\[1pt]}
\begin{minipage}[t]{0.485\textwidth}
\begin{tabularx}{\linewidth}[t]{@{}>{$}l<{$}@{\hspace{0.8em}}>{\raggedright\arraybackslash}X@{}}
\toprule
\nhead{Respondents, groups and items}
i=1,\dots,n & respondent \\
g\in\{0,\dots,m\} & group (ideology, cohort, country, \dots); $g=0$ the reference \\
j;\ Y_{ij}\in\{0,\dots,C_j\} & item; ordinal response of respondent $i$ \\
\mathcal{S},\ \mathcal{I};\ s(j) & stereotype items; Israel-related items; item $j$'s domain \\
\mathcal{A}\subset\mathcal{S}\cup\mathcal{I} & anchor items, assumed invariant across groups \\
\midrule
\nhead{Measurement model}
\bm\theta_i=(\theta_{iG},\theta_{iS},\theta_{iI}) & general factor $G$; specific factors $S$ and $I$ \\
\bm\theta_i\mid g\sim\mathrm{N}_3(\bm\mu_g,\bm\Sigma_g) & latent law in group $g$; $\bm\Sigma_g$ diagonal, $\bm\mu_0=\bm0$, $\bm\Sigma_0=\bm I_3$ \\
a_{jG},\ a_{js} & slopes on $G$ and on item $j$'s specific factor $s=s(j)$ \\
b_{jc} & threshold for category $c$ of item $j$ \\
F & link: logistic $\Lambda$ or normal $\Phi$ (normal ogive) \\
\bm\lambda_j,\ h_j^2 & standardised loadings; communality \\
\bm a^{(g)}_j,\ b^{(g)}_{jc} & group-specific parameters: DIF \\
\zeta_S,\ \zeta_I;\ \alpha_j & correlated domain factors; their loadings \\
\bottomrule
\end{tabularx}
\end{minipage}\hfill
\begin{minipage}[t]{0.485\textwidth}
\begin{tabularx}{\linewidth}[t]{@{}>{$}l<{$}@{\hspace{0.8em}}>{\raggedright\arraybackslash}X@{}}
\toprule
\nhead{Boundary estimands}
\mathrm{ECV}_j,\ \mathrm{ECV}^{(g)}_j & share of item $j$'s common variance due to $G$; in group $g$ \\
\rho_{SI} & latent domain correlation of $\zeta_S$ and $\zeta_I$ \\
\tau;\ \mathcal{B}_\tau & pre-registered threshold; boundary set $\{j\in\mathcal{I}:\mathrm{ECV}_j\ge\tau\}$ \\
\midrule
\nhead{Sensitive questions}
\mathsf{T},\ \mathsf{C},\ \mathsf{D} & list-treatment, list-control and direct arms \\
J & number of control items in a list (\cref{sec:sensitive} only) \\
\pi;\ \pi^D & prevalence of private agreement; of direct agreement \\
\widehat\pi=\overline Y_{\mathsf{T}}-\overline Y_{\mathsf{C}} & list-experiment (difference-in-means) estimator \\
\mathrm{SDB}=\pi-\pi^D & social-desirability bias \\
\pi^H;\ \delta_j & hidden-endorsement rate; threshold shift of item $j$ \\
\midrule
\nhead{Small areas}
k;\ N_k & poststratification cell; its population count \\
r;\ \omega_k=N_k/N_r & area; cell $k$'s share of its area \\
\bar\theta_r & poststratified area mean of the latent traits (\cref{eq:mrp}) \\
\bottomrule
\end{tabularx}
\end{minipage}
\end{table}
}{}%
\paragraph{Contributions}
Our contributions bring established psychometric and survey methods to a
contested question and adapt them to it. Classical results are attributed
where they are used. Most of the formal results specialise known theory to
the two-domain bifactor design and to the boundary estimands, and we claim
novelty for none of them beyond that specialisation and the explicit forms
given.
\begin{enumerate}[label=\textbf{C\arabic*},ref=C\arabic*,leftmargin=2.4em,itemsep=3pt]
  \item\label{contrib:model} \textbf{Measurement model and boundary
  estimands.} A multidimensional graded response model
  \citep{samejima1969} with a bifactor structure
  \citep{gibbons1992,reise2012}: a general factor and specific factors for
  classic stereotypes and for Israel-related attitudes (\cref{sec:model}).
  Within it we define three boundary estimands (\cref{def:boundary}): the
  item-level explained common variance $\mathrm{ECV}_j$ of each
  Israel-related item, overall and by group; the latent correlation
  $\rho_{SI}$ between the stereotype and Israel-attitude factors in a
  correlated-factors model; and the boundary set $\mathcal{B}_\tau$ for a
  pre-registered threshold $\tau$, reported with a three-way classification
  that shows indeterminacy instead of hiding it. We state what the estimands
  would and would not mean (\cref{sec:model-meaning}), and show that the
  correlated-factors model identifies $\rho_{SI}$ but only bounds the
  item-level shares (\cref{prop:model-sl}; \citealp{schmid1957}).
  \item\label{contrib:ident} \textbf{Identification and invariance.} With at
  least three items carrying specific variance in each domain and a general
  loading in each, the item parameters and every $\mathrm{ECV}_j$ are
  identified when, in at least one domain, the items differ in their
  general-factor shares in a tetrad sense; when this fails in both domains,
  an explicit curve of observationally equivalent solutions moves the
  boundary ECVs, and with two Israel-related items their identified sets are
  given in closed form (\cref{thm:ident}; \cref{rem:id-literature} relates
  it to the classical and bifactor literature). An explicit first-order formula for the
  bias in group comparisons of latent means and of the group-specific
  $\mathrm{ECV}^{(g)}_j$ when DIF is ignored, which shows how DIF can
  manufacture or mask a difference in the general factor and how
  non-uniform DIF in some items moves the reported shares of all items
  (\cref{thm:dif-bias}). Identification of group comparisons with anchor
  items under partial invariance \citep{byrne1989}, the undetectable
  failure mode when the anchors themselves have DIF, and a sensitivity bound
  in terms of the anchors' leverage (\cref{prop:anchor}); anchors drawn
  from the stereotype items alone suffice for comparisons of the general
  factor and of the boundary (\cref{rem:id-anchor-S}). Tests and confidence
  sets for $\mathcal{B}_\tau$ that justify the three-way classification
  (\cref{prop:id-boundary-test}).
  \item\label{contrib:sdb} \textbf{Social-desirability correction.} The
  exact variance of the list experiment \citep{miller1984,glynn2013} and its
  cost relative to direct questioning, which no choice of control items can
  remove without revealing individual answers (\cref{prop:list-var}). A
  joint model of direct and list responses that identifies item-level
  social-desirability bias in split and paired designs, with and without
  the monotonicity assumption of \citet{aronow2015}, and decomposes it into
  hidden and false endorsement where the design allows
  (\cref{thm:joint-sdb}). The budget-optimal allocation of respondents across
  direct and list arms and the choice of control list under floor and
  ceiling effects \citep{blair2012}: larger budgets favour lists that are
  noisier but less exposed (\cref{thm:design}). A link to the measurement
  model that treats social desirability as a shift of an item's thresholds
  and identifies the shift from a list experiment (\cref{prop:list-irt}).
  \item\label{contrib:mrp} \textbf{Small-area estimation and indices.}
  Multilevel regression and poststratification \citep{gelman1997,park2004}
  on the latent traits, fitted jointly with the measurement model, with an
  explicit poststratification identity for area means (\cref{eq:mrp}) that
  extends to area covariances and area-level boundary summaries
  (\cref{sec:sa-mrp}). Two-step scoring attenuates area differences by an
  explicit factor (\cref{prop:eap}). For stereotype-count indices, a
  sufficient condition under which group differences in the index have the
  sign of the latent difference, and constructions in which unequal
  dispersion, uniform DIF or non-uniform DIF reverse it
  (\cref{prop:index}). We make no claim about the results of any real
  index.
  \item\label{contrib:sim} \textbf{A simulation study.} In synthetic surveys
  with known parameters we assess recovery of the boundary estimands, the
  bias from ignoring DIF, the failure caused by a misspecified anchor, the
  accuracy of the list-experiment correction and of the optimal design, the
  accuracy of small-area estimates, and how often a stereotype-count index
  ranks two groups in the wrong order when its items carry DIF
  (\cref{sec:simulation}). The fits use a limited-information estimator,
  not the full-information estimator of the plan, and the small-area
  experiment starts from a score of the general factor rather than from item
  responses. The results agree with the theory in the designs studied. They
  also show that convergence and agreement across starting values do not
  reveal an unidentified decomposition, and they exposed a flaw in our own
  draft plan: a DIF screen that also required a minimum expected-score
  difference missed non-uniform DIF. The plan was revised accordingly.
  \item\label{contrib:design} \textbf{A pre-registered survey design and
  analysis plan.} An item bank described by construct, with no claim of
  validation; the list-experiment module; the sampling frame and the
  poststrata; the estimands, the threshold $\tau$ and decision rules fixed
  before fieldwork; and an ethics protocol (\cref{sec:design,sec:ethics}).
\end{enumerate}

\paragraph{What the method cannot settle}
Two limits are built into the approach, and we state them at the outset.
First, \emph{the method does not answer normative questions.} An estimate of
$\mathrm{ECV}_j$ says how responses to a statement covary with responses to
classic stereotypes in a surveyed population. Whether speech with that
content should be called antisemitic, restricted, disciplined or protected is
a question of values, law and context. Evidence about covariance can inform
that question, but it cannot answer it, and a high or low $\mathrm{ECV}_j$ is
compatible with more than one normative position. Second, \emph{the method
does not diagnose individuals.} The estimands describe populations. An item
in $\mathcal{B}_\tau$ does not show that a particular person who endorses it is hostile
toward Jews, or that a particular utterance with similar content is
antisemitic in its context; an item outside $\mathcal{B}_\tau$ is not thereby shown to
be acceptable in every context. Definitions classify speech and conduct in
context, often with reference to intent (\cref{sec:bg-definitions}); our
estimands do neither, and \cref{sec:ethics-individual} turns these limits
into commitments about how results are reported.

\paragraph{Related work}
The graded response model goes back to \citet{samejima1969}, building on
the logistic models of \citet{birnbaum1968} and the item response theory of
\citet{lord1980}. Full-information bifactor analysis is due to
\citet{gibbons1992}, its revival in the measurement of psychological
constructs to \citet{reise2012}, and the bifactor indices used here,
including explained common variance, are reviewed by \citet{rodriguez2016}.
The theory of measurement invariance is that of \citet{meredith1993} and
\citet{millsap2011}; DIF is tested with the Mantel--Haenszel procedure
\citep{holland1988} and with likelihood-ratio tests in an item response model
\citep{thissen1993}; partial invariance is due to \citet{byrne1989}. The
list experiment originates with \citet{miller1984} and is analysed by
\citet{glynn2013}, \citet{imai2011} and \citet{blair2012}; randomised
response is due to \citet{warner1965}. Multilevel regression and
poststratification was introduced by \citet{gelman1997} and
\citet{park2004}. On the substantive side, \cref{sec:background} reviews the
three definitions, the main survey programmes and prior evidence on the
relation between anti-Israel and antisemitic attitudes; several studies
report positive associations between the two domains
\citep{kaplan2006,staetsky2017,beattie2017}. This paper is the third in a
series. The first
estimates the number of antisemitic hate crimes that no list records
\citep{kornsteinlab2026darkfigure}; the second models surges in recorded
incidents after salient external events \citep{kornsteinlab2026contagion}.
Both study incidents. This one studies attitudes, which no incident list
records.

\paragraph{Roadmap}
\Cref{sec:background} reviews the definitions, the surveys and prior
evidence. \Cref{sec:model} defines the measurement model and the boundary
estimands. \Cref{sec:identification} develops identification and invariance.
\Cref{sec:sensitive} treats social desirability and list experiments, and
\cref{sec:smallarea} small-area estimation and count indices.
\Cref{sec:simulation} reports the simulation study, \cref{sec:design} the
pre-registered design and analysis plan, and \cref{sec:ethics} the ethics
protocol. \Cref{sec:discussion} discusses what a boundary map could and could
not show, limitations and extensions. Proofs are in
\cref{app:proofs-measurement,app:proofs-design} and simulation details in
\cref{app:simdetails}.
